% Worst-case complexity of classical APPR.
% Primary source: Andersen, Chung, and Lang, Internet Mathematics 4(1):35--64,
% 2007. Verified against the source PDF, especially Definition 3.1 (p. 41),
% Definition 3.3 and Lemma 3.4 (p. 42), Algorithm 1 and Theorem 3.2 with
% Equation (3.3) (pp. 41--43).
%
% Notation: this section continues the plain-italic, column-vector convention
% of \cref{sec:problem-formulation}. The source states its algorithm with row
% vectors and the lazy walk matrix acting on the right; we transpose it.

We prove that the classical lazy Andersen--Chung--Lang approximate
personalized PageRank algorithm (APPR) of \citet{andersen2007using} has
worst-case adjacency-list work
\(\Theta(1/(\alpha\epsappr))\).
The upper bound is the residual-mass argument of that paper. For the matching
lower bound, a star whose center is the seed is sufficient. Moreover, the
lower bound holds for every legal ordering of active vertices; it is not an
artifact of FIFO, priority-queue, or any other tie-breaking rule. This
settles the worst-case complexity of the baseline against which accelerated
and hybrid local solvers are measured in this manuscript.

\subsection{Algorithm and work model}
\label{subsec:appr-work-model}

Use the common graph and seed conventions of
\cref{subsec:shared-source-aligned-problem}, specialized to the single seed
\(s=e_v\). Fix \(\alpha\in(0,1]\) and an APPR residual tolerance
\(\epsappr>0\), and write
\begin{equation}
    a:=\frac{1-\alpha}{2},
    \label{eq:appr-lazy-factor}
\end{equation}
for the lazy half-step factor. APPR maintains nonnegative vectors
\(p_t,r_t\in\R_{\geq0}^{n}\), initialized by
\begin{equation}
    p_0=0,
    \qquad
    r_0=s=e_v.
    \label{eq:appr-initialization}
\end{equation}
A vertex \(u\) is \emph{active} at time \(t\) if
\begin{equation}
    r_t(u)\geq\epsappr\,d_u.
    \label{eq:appr-active}
\end{equation}
Given any active vertex \(u_t\), set \(\xi_t:=r_t(u_t)\). The operation
\(\push(u_t)\) makes the updates
\begin{align}
    p_{t+1}(u_t)
        &=p_t(u_t)+\alpha \xi_t,
        \label{eq:appr-push-p}\\
    r_{t+1}(u_t)
        &=a\,\xi_t,
        \label{eq:appr-push-self}\\
    r_{t+1}(w)
        &=r_t(w)+\frac{a\,\xi_t}{d_{u_t}}
        \qquad\text{for every }w\sim u_t,
        \label{eq:appr-push-neighbor}
\end{align}
while all other coordinates remain unchanged. These are exactly the updates of
\citet[Definition~3.3]{andersen2007using}. The algorithm terminates at the
first time \(T\) with no active vertex, that is, when
\begin{equation}
    r_T(u)<\epsappr\,d_u
    \qquad\text{for every }u\in V.
    \label{eq:appr-termination}
\end{equation}
The source imposes no further requirement on the choice of \(u_t\): its
Algorithm~1 picks \emph{any} active vertex. We call any sequence obeying
\cref{eq:appr-active,eq:appr-termination} a \emph{legal ordering} and denote it
by \(\sigma\).

We measure adjacency-list work by
\begin{equation}
    W:=\sum_{t=0}^{T-1}d_{u_t}.
    \label{eq:appr-work}
\end{equation}
A push at \(u\) costs \(d_u\), because it scans the adjacency list of \(u\).
This is the same unit of work as the degree-weighted measure
\cref{eq:source-degree-work}: both charge \(d_i\) whenever the neighborhood of
vertex \(i\) is scanned. They differ only in which coordinates are scanned per
iteration, since APPR touches one active coordinate at a time whereas a
proximal-gradient step scans the whole current support.

\subsection{What APPR computes}
\label{subsec:appr-target}

Let \(P:=AD^{-1}\) be the column-stochastic random-walk matrix. The shared
lazy personalized PageRank vector \(\pi\) from
\cref{eq:shared-ppr-solution} satisfies
\begin{equation}
    \pi
    =\alpha s+(1-\alpha)\frac{I+P}{2}\pi.
    \label{eq:lazy-ppr-system}
\end{equation}
More generally, write \(\pr(z):=\alpha\bigl(I-(1-\alpha)(I+P)/2\bigr)^{-1}z\)
for the lazy PageRank vector of an arbitrary nonnegative \(z\), so that
\(\pi=\pr(s)\). The operator \(\pr\) is linear, preserves nonnegativity, and
preserves \(\ell_1\) mass, because \((I+P)/2\) is column-stochastic.

\begin{lemma}[Push invariant, {\citealp[Lemma~3.4]{andersen2007using}}]
\label{lem:appr-invariant}
Every APPR execution satisfies \(p_t+\pr(r_t)=\pi\) for all \(t\), and hence
\begin{equation}
    \norm{\pi-p_T}_1
    =\norm{r_T}_1
    <\epsappr\vol(V).
    \label{eq:appr-accuracy}
\end{equation}
\end{lemma}

\begin{proof}
The invariant is the column-vector transpose of
\citet[Lemma~3.4]{andersen2007using}, whose statement
\(p=\pr(s-r)\) is equivalent to \(p+\pr(r)=\pr(s)\) by linearity of \(\pr\).
The identity \(\norm{\pi-p_T}_1=\norm{\pr(r_T)}_1=\norm{r_T}_1\) uses
nonnegativity and mass preservation, and the final inequality sums
\cref{eq:appr-termination} over \(V\).
\end{proof}

The next proposition identifies the fixed-point form used by APPR with the
shared optimization target.

\begin{proposition}[APPR targets the unregularized RPPR solution]
\label{prop:appr-rppr-bridge}
The solution of \cref{eq:lazy-ppr-system} is the vector \(\pi\) in
\cref{eq:shared-ppr-solution}. Equivalently,
\(x^0=D^{-1/2}\pi\) is the unique minimizer of the shared smooth
objective \(f\), equivalently of the \(\rho=0\) extension of \(F_\rho\).
\end{proposition}

\begin{proof}
Rearranging \cref{eq:lazy-ppr-system} gives
\(\frac{1+\alpha}{2}\pi-\frac{1-\alpha}{2}P\pi=\alpha s\).
Since \(D^{-1/2}AD^{-1/2}D^{-1/2}\pi=D^{-1/2}AD^{-1}\pi=D^{-1/2}P\pi\),
\cref{eq:shared-pagerank-matrices} yields
\[
    QD^{-1/2}\pi
    =\frac{1+\alpha}{2}D^{-1/2}\pi
     -\frac{1-\alpha}{2}D^{-1/2}P\pi
    =D^{-1/2}\left(
        \frac{1+\alpha}{2}\pi-\frac{1-\alpha}{2}P\pi
    \right)
    =\alpha D^{-1/2}s.
\]
Because \(\alpha I\preceq Q\preceq I\), this identifies the fixed-point
solution with the unique point \(x^0=Q^{-1}b\) already defined in
\cref{eq:shared-ppr-solution}.
\end{proof}

\begin{remark}[Scope of \(\epsappr\)]
\label{rem:appr-epsilon-scope}
The tolerance \(\epsappr\) is a degree-normalized residual threshold in the
sense of \cref{eq:appr-active}. It is not identified in this manuscript with
the objective-gap target \(\epsobj\) or the proximal fixed-point tolerance
\(\epspg\) of
\cref{subsec:source-work-and-stopping}. \Cref{eq:appr-accuracy} is the only
accuracy statement attached to \(\epsappr\) here. Similarly, \(\epsappr\) and
the sparsity parameter \(\rho\) play analogous locality roles, since
\(\vol(\supp(x^\star(\rho)))\leq1/\rho\) while
\(\vol(\supp(p_T))\leq2/((1-\alpha)\epsappr)\) by
\citet[Theorem~3.2]{andersen2007using} when \(0<\alpha<1\), but we do not
assert a translation between them.
\end{remark}

\subsection{Mass identity and the classical upper bound}
\label{subsec:appr-upper-bound}

\begin{lemma}[Residual-mass identity]
\label{lem:appr-mass-identity}
For every APPR execution and every \(T\geq0\),
\begin{equation}
    \norm{p_T}_1
    =\alpha\sum_{t=0}^{T-1}\xi_t,
    \qquad
    \norm{r_T}_1
    =1-\alpha\sum_{t=0}^{T-1}\xi_t,
    \label{eq:appr-mass-identities}
\end{equation}
and hence
\begin{equation}
    \norm{p_T}_1+\norm{r_T}_1=1.
    \label{eq:appr-total-mass}
\end{equation}
In particular, \(\norm{r_t}_1\leq1\) at every time \(t\).
\end{lemma}

\begin{proof}
Consider a push whose residual is \(\xi_t\). \Cref{eq:appr-push-p}
increases \(\norm{p_t}_1\) by \(\alpha \xi_t\). The pushed residual coordinate
retains \(a\,\xi_t\), while its neighbors receive total residual
\(d_{u_t}\cdot a\,\xi_t/d_{u_t}=a\,\xi_t\). Consequently,
\begin{equation}
    \norm{r_{t+1}}_1
    =\norm{r_t}_1-\xi_t+2a\,\xi_t
    =\norm{r_t}_1-\alpha \xi_t,
    \label{eq:appr-one-step-mass}
\end{equation}
where the last equality uses \(2a=1-\alpha\). Summing over all pushes and
using \cref{eq:appr-initialization} proves
\cref{eq:appr-mass-identities,eq:appr-total-mass}.
\end{proof}

\begin{proposition}[Classical APPR upper bound,
{\citealp[Theorem~3.2]{andersen2007using}}]
\label{prop:appr-upper-bound}
Every terminating APPR execution satisfies
\begin{equation}
    W\leq\frac{1}{\alpha\epsappr}.
    \label{eq:appr-upper-bound}
\end{equation}
This bound is independent of the ordering in which active vertices are
selected.
\end{proposition}

\begin{proof}
Because \(u_t\) is active, \(\xi_t\geq\epsappr d_{u_t}\), and therefore
\begin{equation}
    W
    =\sum_{t=0}^{T-1}d_{u_t}
    \leq\frac{1}{\epsappr}\sum_{t=0}^{T-1}\xi_t
    =\frac{\norm{p_T}_1}{\alpha\epsappr}
    \leq\frac{1}{\alpha\epsappr},
    \label{eq:appr-upper-bound-proof}
\end{equation}
where \cref{lem:appr-mass-identity} gives the equality and the final
inequality. Notice also that every push has
\(\xi_t\geq\epsappr d_{u_t}\geq\epsappr\). Since
\(\sum_t\xi_t\leq1/\alpha\), every legal execution contains at most
\(1/(\alpha\epsappr)\) pushes and therefore terminates. This is
\citet[Equation~(3.3)]{andersen2007using}, restated in the work model
\cref{eq:appr-work}.
\end{proof}

\subsection{A matching star lower bound}
\label{subsec:appr-star-lower-bound}

Let \(K_{1,m}\) be the star with center \(c\) and leaf set \(L\), where
\begin{equation}
    m:=\floor{\frac{1}{8\epsappr}},
    \qquad
    v:=c.
    \label{eq:hard-star-size}
\end{equation}
This is the center-seeded star of \cref{subsec:source-auxiliary-notation}, with
the leaf count tied to \(\epsappr\). For \(0<\epsappr\leq1/16\), the choice in
\cref{eq:hard-star-size} satisfies
\begin{equation}
    \frac{1}{16\epsappr}\leq m\leq\frac{1}{8\epsappr},
    \qquad
    \vol(V)=2m.
    \label{eq:hard-star-size-bounds}
\end{equation}
Indeed, \(1/(8\epsappr)\geq2\), and hence
\(\floor{1/(8\epsappr)}\geq1/(16\epsappr)\).

\begin{theorem}[Ordering-independent star lower bound]
\label{thm:appr-star-lower-bound}
Fix \(\alpha\in(0,1]\) and \(\epsappr\in(0,1/16]\). Run lazy APPR on the
center-seeded star \cref{eq:hard-star-size}. For every legal ordering of
active vertices, the resulting adjacency-list work obeys
\begin{equation}
    W>\frac{3}{128\,\alpha\epsappr}.
    \label{eq:appr-star-lower-bound}
\end{equation}
Together with \cref{prop:appr-upper-bound}, this yields
\begin{equation}
    \frac{3}{128\,\alpha\epsappr}
    <W
    \leq\frac{1}{\alpha\epsappr}.
    \label{eq:appr-star-two-sided-bound}
\end{equation}
More formally, define the worst-case work over all finite graphs without
isolated vertices, all seeds, and all legal orderings by
\begin{equation}
    W_{\mathrm{appr}}^{\mathrm{worst}}(\alpha,\epsappr)
    :=\sup_{G,v,\sigma}W(G,v;\alpha,\epsappr,\sigma).
    \label{eq:appr-worst-work-definition}
\end{equation}
Then classical APPR has worst-case work
\begin{equation}
    W_{\mathrm{appr}}^{\mathrm{worst}}(\alpha,\epsappr)
    =\Theta\!\left(\frac{1}{\alpha\epsappr}\right).
    \label{eq:appr-tight-complexity}
\end{equation}
\end{theorem}

\begin{proof}
Let
\begin{equation}
    Z_c:=\sum_{t:u_t=c}\xi_t,
    \qquad
    Z_L:=\sum_{t:u_t\in L}\xi_t,
    \qquad
    Z:=Z_c+Z_L
    \label{eq:appr-cumulative-masses}
\end{equation}
be the cumulative pushed residual at the center, at the leaves, and in total.
The seed is active initially, since \(r_0(c)=1\) and
\(\epsappr d_c=\epsappr m\leq1/8\), so the execution is nonempty. We establish
the result in three steps.

\paragraph{Step 1: APPR must process \(\Omega(1/\alpha)\) residual mass.}
At termination, \cref{eq:appr-termination,eq:hard-star-size-bounds} give
\begin{equation}
    \norm{r_T}_1
    <\epsappr\sum_{u\in V}d_u
    =\epsappr\vol(V)
    =2\epsappr m
    \leq\frac14.
    \label{eq:appr-star-terminal-residual}
\end{equation}
\Cref{lem:appr-mass-identity} therefore implies
\begin{equation}
    \norm{p_T}_1
    =1-\norm{r_T}_1
    >\frac34.
    \label{eq:appr-star-settled-mass}
\end{equation}
Because \(\norm{p_T}_1=\alpha Z\), we obtain
\begin{equation}
    Z>\frac{3}{4\alpha}.
    \label{eq:appr-star-total-pushed-mass}
\end{equation}
This is the source of the factor \(1/\alpha\).

\paragraph{Step 2: A constant fraction of the pushed mass passes through the
center.}
Define the aggregate residual on the leaves by
\begin{equation}
    \res_L(t):=\sum_{\ell\in L}r_t(\ell).
    \label{eq:appr-leaf-residual}
\end{equation}
A center push of size \(\xi\) sends total residual \(a\,\xi\) to the leaves and
hence changes \(\res_L\) by \(+a\,\xi\). A leaf push of size \(\xi\) changes the
residual at that leaf from \(\xi\) to \(a\,\xi\) and sends \(a\,\xi\) to the center,
so it changes \(\res_L\) by \(-(1-a)\xi\). Since \(\res_L(0)=0\), summing these
exact changes over the entire execution yields
\begin{equation}
    \res_L(T)=a\,Z_c-(1-a)Z_L.
    \label{eq:appr-leaf-flow-identity}
\end{equation}
The nonnegativity of the terminal residual gives
\begin{equation}
    Z_L\leq\frac{a}{1-a}Z_c.
    \label{eq:appr-leaf-center-comparison}
\end{equation}
Consequently,
\begin{equation}
    Z=Z_c+Z_L
    \leq\frac{Z_c}{1-a},
    \qquad\text{so}\qquad
    Z_c\geq(1-a)Z.
    \label{eq:appr-center-fraction}
\end{equation}
Since \(1-a=(1+\alpha)/2\geq1/2\),
\cref{eq:appr-star-total-pushed-mass,eq:appr-center-fraction} imply
\begin{equation}
    Z_c>\frac{3}{8\alpha}.
    \label{eq:appr-center-pushed-mass}
\end{equation}

\paragraph{Step 3: Convert center mass into adjacency-list work.}
Let \(k_c\) denote the number of center pushes. By
\cref{lem:appr-mass-identity}, every center-push residual satisfies
\begin{equation}
    \xi_t=r_t(c)\leq\norm{r_t}_1\leq1.
    \label{eq:appr-center-push-at-most-one}
\end{equation}
It follows that
\begin{equation}
    k_c\geq Z_c>\frac{3}{8\alpha}.
    \label{eq:appr-number-center-pushes}
\end{equation}
Each center push scans all \(m\) incident edges. Therefore,
\begin{equation}
    W
    \geq m\,k_c
    >\frac{1}{16\epsappr}\cdot\frac{3}{8\alpha}
    =\frac{3}{128\,\alpha\epsappr},
    \label{eq:appr-final-star-calculation}
\end{equation}
which proves \cref{eq:appr-star-lower-bound}. The matching upper bound in
\cref{prop:appr-upper-bound} holds for every instance and every ordering, so
\cref{eq:appr-tight-complexity} follows. None of the three steps uses the order
in which active vertices are selected, proving the claimed ordering
independence.
\end{proof}

\begin{corollary}[Properties of the hard instance]
\label{cor:appr-hard-instance-properties}
The lower bound in \cref{thm:appr-star-lower-bound} is witnessed by an
unweighted, undirected, single-seed instance and measures actual
adjacency-update work rather than merely the number of push operations. It
holds under FIFO, LIFO, priority-based, and arbitrary legal active-vertex
orderings.
\end{corollary}

\begin{proof}
The graph in \cref{eq:hard-star-size} has all the stated structural
properties. The proof of \cref{thm:appr-star-lower-bound} lower-bounds
\(\sum_td_{u_t}\) directly and is independent of the active-vertex ordering.
\end{proof}

\begin{remark}[Constants]
\label{rem:appr-constants}
No attempt is made to optimize the constant \(3/128\). The leaf count
\(m=\floor{1/(8\epsappr)}\) was chosen so that the terminal residual budget
\(\epsappr\vol(V)\) is at most \(1/4\); other choices trade the constants in
\cref{eq:appr-star-terminal-residual,eq:appr-final-star-calculation} against
each other. Only the \(\Theta(1/(\alpha\epsappr))\) scaling in
\cref{eq:appr-tight-complexity} is claimed to be tight.
\end{remark}

\subsection{Why the star is the clean witness}
\label{subsec:appr-star-discussion}

The construction isolates the two mechanisms in the APPR upper bound.
\begin{enumerate}[leftmargin=*]
    \item The stopping condition on a graph of volume
    \(\Theta(1/\epsappr)\) forces APPR to settle a constant amount of
    probability mass. Since a push settles only an \(\alpha\)-fraction of its
    pushed residual, the cumulative pushed residual is \(\Omega(1/\alpha)\).
    \item The leaf-flow identity \cref{eq:appr-leaf-flow-identity} forces a
    constant fraction of that cumulative residual through the center. The
    center has degree \(\Theta(1/\epsappr)\), so repeatedly processing it
    incurs \(\Omega(1/(\alpha\epsappr))\) adjacency-list work.
\end{enumerate}
Thus the star simultaneously realizes the \(1/\alpha\) mass-processing factor
and the \(1/\epsappr\) high-degree work factor.

\subsection{Paths and spiders}
\label{subsec:appr-paths-and-spiders}

\Cref{thm:appr-star-lower-bound} does not require a path or a genuinely long
spider. A depth-one spider is exactly a star and hence already supplies the
matching lower bound. Long spiders may be useful for lower bounds that must
also expose propagation distance, repeated wavefront motion, or local spectral
effects, as occurs in analyses of successive over-relaxation and accelerated
local methods; those properties require a separate argument and are not needed
for classical APPR.

The contrast with a path can be seen from the exact lazy PageRank kernel on
the infinite path.

\begin{lemma}[Lazy PageRank kernel on the infinite path]
\label{lem:lazy-path-kernel}
Let the graph be the infinite path \(\sZ\), let the seed be the origin, and let
\(\alpha\in(0,1)\). The solution of \cref{eq:lazy-ppr-system} is
\begin{equation}
    \pi(j)
    =\sqrt{\alpha}
     \left(\frac{1-\sqrt{\alpha}}{1+\sqrt{\alpha}}\right)^{\abs{j}},
    \qquad j\in\sZ.
    \label{eq:lazy-path-kernel}
\end{equation}
For \(\alpha\downarrow0\), its decay length is \(\Theta(1/\sqrt{\alpha})\).
\end{lemma}

\begin{proof}
Write \(\theta:=\sqrt{\alpha}\). Away from the seed, every neighbor has degree
two, so \cref{eq:lazy-ppr-system} reads
\begin{equation}
    \pi(j)
    =(1-\alpha)
     \left(
        \frac12\pi(j)
        +\frac14\pi(j-1)
        +\frac14\pi(j+1)
     \right),
    \qquad j\neq0.
    \label{eq:lazy-path-recurrence}
\end{equation}
The characteristic equation has a unique root in \([0,1)\), namely
\begin{equation}
    \varrho
    =\frac{1-\theta}{1+\theta}.
    \label{eq:lazy-path-rho}
\end{equation}
Symmetry and summability therefore give \(\pi(j)=C\varrho^{\abs{j}}\).
Normalizing the total mass to one gives
\begin{equation}
    1
    =C\left(1+2\sum_{j=1}^{\infty}\varrho^{\,j}\right)
    =C\frac{1+\varrho}{1-\varrho},
    \qquad
    C=\frac{1-\varrho}{1+\varrho}=\theta.
    \label{eq:lazy-path-normalization}
\end{equation}
The seed equation is also satisfied, since
\begin{equation}
    \alpha+(1-\alpha)C\frac{1+\varrho}{2}
    =\theta^2+(1-\theta^2)\frac{\theta}{1+\theta}
    =\theta
    =C.
    \label{eq:lazy-path-source-check}
\end{equation}
This proves \cref{eq:lazy-path-kernel}. Finally,
\begin{equation}
    -\log\varrho
    =2\operatorname{arctanh}(\sqrt{\alpha})
    =2\sqrt{\alpha}+\bigO(\alpha^{3/2}),
    \label{eq:lazy-path-decay-rate}
\end{equation}
so the reciprocal decay rate is \(\Theta(1/\sqrt{\alpha})\).
\end{proof}

\begin{remark}[What the path calculation does and does not show]
\label{rem:appr-path-caveat}
A path has bounded degree, and its intrinsic lazy PageRank propagation scale is
\(1/\sqrt{\alpha}\) by \cref{lem:lazy-path-kernel}. It therefore does not
expose the two factors in the APPR bound as directly as the star, whose
expensive center has degree \(\Theta(1/\epsappr)\). This observation is not an
impossibility theorem for every path parameter regime: a finite chain can still
be tight under additional relations among its length, \(\alpha\), and
\(\epsappr\). For example, regimes with length \(L=\Theta(1/\epsappr)\) and
\(\alpha L^2=\bigO(1)\) merit a separate schedule-sensitive analysis. The point
here is only that no such additional argument or parameter coupling is needed
for the center-seeded star.
\end{remark}

\paragraph{Summary.}
The center-seeded star proves, by itself and without assumptions on the active
ordering, that the classical bound \(\bigO(1/(\alpha\epsappr))\) for lazy APPR
is worst-case tight. The obstruction is exact: APPR must process
\(\Theta(1/\alpha)\) cumulative residual mass, and the star forces a constant
fraction of this mass to be processed repeatedly at a vertex of degree
\(\Theta(1/\epsappr)\). Any local solver that improves on
\cref{eq:appr-tight-complexity} must therefore break at least one of these two
mechanisms, which is the design target of the hybrid method developed in
\cref{sec:algorithm}.
